\chapter{Expectation and Moment Inequalities}

\section{Expectation and Variance}

\begin{definitionbox}
    Let $X$ be a random variable.
    The \emph{expectation} of $X$ is
    $$\EE(X) = \int_{\RR} xf(x)\, dx.$$
    The \emph{variance} of $X$ is
    $$\Var(X) = \EE(X - \EE(X))^2 = \EE(X^2) - (\EE(X))^2.$$
\end{definitionbox}

\section{Covariance and Random Vectors}

\begin{definitionbox}[Covariance]
    Let $X$ and $Y$ be two random variables. The \emph{covariance} of $X$ and $Y$ is
    $$\Cov(X,Y) = \EE\left((X - \EE(X))(Y - \EE(Y))\right) = \EE(XY) - \EE(X)\EE(Y).$$

    Let $X_1, \ldots, X_n$ be random variables. The \emph{covariance matrix} of $(X_1, \ldots, X_n)$ is
    $$\Sigma = \begin{pmatrix}
        \Cov(X_1, X_1) & \Cov(X_1, X_2) & \cdots & \Cov(X_1, X_n) \\
        \Cov(X_2, X_1) & \Cov(X_2, X_2) & \cdots & \Cov(X_2, X_n) \\
        \vdots & \vdots & \ddots & \vdots \\
        \Cov(X_n, X_1) & \Cov(X_n, X_2) & \cdots & \Cov(X_n, X_n)
    \end{pmatrix}.$$
\end{definitionbox}

\begin{propositionbox}
    Let $X_1, \ldots, X_n$ be random variables. Then the covariance matrix $\Sigma$ is symmetric and positive semi-definite.
\end{propositionbox}
\begin{proofbox}
    Since $\Cov(X_i, X_j) = \Cov(X_j, X_i)$, we only need to show that the covariance matrix is positive semi-definite.
    Let $a = (a_1, \ldots, a_n)^T \in \RR^n$. Then
    $$\begin{aligned}
        a^T \Sigma a &= \sum_{i,j=1}^{n} a_i a_j \Cov(X_i, X_j) \\
        &= \sum_{i,j=1}^{n} a_i a_j \EE\left((X_i - \EE(X_i))(X_j - \EE(X_j))\right) \\
        &= \EE\left(\sum_{i=1}^{n} a_i (X_i - \EE(X_i))\right)^2 \\
        &\geq 0.
    \end{aligned}$$
\end{proofbox}

\begin{propositionbox}
    If $X \perp Y$, then
    $$\Var(X + Y) = \Var(X) + \Var(Y).$$
\end{propositionbox}
\begin{proofbox}
    $$\begin{aligned}
        \Var(X + Y) &= \EE\left((X - \mu_1) + (Y - \mu_2)\right)^2 \\
        & = \EE(X - \mu_1)^2 + \EE(Y - \mu_2)^2 + 2\EE[(X - \mu_1)(Y - \mu_2)] \\
        & = \Var(X) + \Var(Y) + 2\EE(X - \mu_1)\EE(Y - \mu_2)\\
        & = \Var(X) + \Var(Y).
    \end{aligned}$$
\end{proofbox}

\subsection{Multivariate Normal Distribution}

\begin{definitionbox}[Multivariate normal distribution]
    Let $X_1, \ldots, X_n$ be random variables. The vector
    $(X_1, \ldots, X_n)$ is said to have a \emph{multivariate normal
    distribution} if its probability density function is given by
    $$f(x_1, \ldots, x_n) = \frac{1}{(2\pi)^{n/2}|\det \Sigma|^{1/2}} \exp\left(-\frac{1}{2}(x - \mu)^T \Sigma^{-1} (x - \mu)\right),$$
    where $\mu = (\mu_1, \ldots, \mu_n)^T$ is the mean vector and $\Sigma$ is the covariance matrix.

    We denote this as $(X_1, \ldots, X_n) \sim \mathcal{N}(\mu, \Sigma)$.

    In particular, if $X_1, \ldots, X_n$ are pairwise independent, then
    $\Sigma$ is diagonal and the density function simplifies to
    $$f(x_1, \ldots, x_n) = \prod_{i=1}^{n} \frac{1}{\sqrt{2\pi}\sigma_i} \exp\left(-\frac{(x_i - \mu_i)^2}{2\sigma_i^2}\right),$$
    where $\sigma_i^2 = \Var(X_i)$.
\end{definitionbox}

\section{Moment Inequalities}

\subsection{Jensen and H\"older Inequalities}

\begin{remarkbox}
    By Jensen's inequality, if $J: \RR \rightarrow \RR$ is convex, then
    $$J(\EE(X)) \leq \EE(J(X)).$$
    In particular, if $p \ge 1$, then
    $$\EE(X^p) \geq \left(\EE(X)\right)^p.$$
\end{remarkbox}

\begin{theorembox}[H\"older]
    If $p,q \in [1,+\infty]$ with $\frac{1}{p} + \frac{1}{q} = 1$, then
    $$\EE(|XY|) \leq \left(\EE(|X|^p)\right)^{\frac{1}{p}} \left(\EE(|Y|^q)\right)^{\frac{1}{q}}.$$
\end{theorembox}

\subsection{Minkowski Inequality}

\begin{theorembox}[Minkowski]
    Let $(X,\mathcal{X},\mu)$ and $(T,\mathcal{T},\nu)$ be measure spaces, and let
    $f : T \times X \to \RR$ be measurable. If $1 \leq p \leq \infty$, then
    $$\left(\int_{X} \left(\int_{T} |f(t,x)| \nu(\dd t)\right)^p \mu(\dd x)\right)^{1/p}
    \leq \int_{T} \left(\int_{X} |f(t,x)|^p \mu(\dd x)\right)^{1/p} \nu(\dd t),$$
    where the usual essential supremum interpretation is used when $p=\infty$.
    Here, the \emph{essential supremum} of a measurable function $g$ is
    defined by
    $$\operatorname*{ess\,sup}_{x\in X}|g(x)|
    =\inf\left\{M\geq 0:
    \mu\left(\left\{x\in X:|g(x)|>M\right\}\right)=0\right\}.$$
    Thus, the essential supremum is the smallest upper bound that holds
    $\mu$-almost everywhere.
\end{theorembox}
\begin{proofbox}
    Write
    $$H(x) = \int_{T} |f(t,x)|\nu(\dd t).$$
    The case $p=1$ follows immediately from Tonelli's theorem:
    $$\begin{aligned}
        \int_{X} H(x)\mu(\dd x)
        &= \int_{X}\int_{T}|f(t,x)|\nu(\dd t)\mu(\dd x) \\
        &= \int_{T}\int_{X}|f(t,x)|\mu(\dd x)\nu(\dd t).
    \end{aligned}$$

    Suppose now that $1<p<\infty$, and let $q$ be the conjugate exponent,
    so that $\frac{1}{p}+\frac{1}{q}=1$. If $\|H\|_{L^p(\mu)}=0$, the
    desired inequality is trivial. Otherwise, Tonelli's theorem and
    H\"older's inequality give
    $$\begin{aligned}
        \|H\|_{L^p(\mu)}^p
        &= \int_{X} H(x)^{p-1}H(x)\mu(\dd x) \\
        &= \int_{T}\int_{X}|f(t,x)|H(x)^{p-1}\mu(\dd x)\nu(\dd t) \\
        &\leq \int_{T}
        \left(\int_{X}|f(t,x)|^p\mu(\dd x)\right)^{1/p}
        \left(\int_{X}H(x)^{(p-1)q}\mu(\dd x)\right)^{1/q}
        \nu(\dd t) \\
        &= \|H\|_{L^p(\mu)}
        \int_{T}\left(\int_{X}|f(t,x)|^p\mu(\dd x)\right)^{1/p}\nu(\dd t).
    \end{aligned}$$
    Dividing by $\|H\|_{L^p(\mu)}$ proves the result for $1<p<\infty$.

    Finally, when $p=\infty$, for $\nu$-almost every $t$,
    $$|f(t,x)|\leq \operatorname*{ess\,sup}_{y\in X}|f(t,y)|$$
    for $\mu$-almost every $x$. Hence
    $$H(x)\leq \int_{T}\operatorname*{ess\,sup}_{y\in X}|f(t,y)|\nu(\dd t),$$
    and taking the essential supremum over $x$ yields the desired inequality.
\end{proofbox}

\subsection{Markov and Chebyshev Inequalities}

\begin{propositionbox}
    Let $X$ be a random variable.
    Then for any $p \ge 1$ and $\delta > 0$, we have
    $$\PP(|X| > \delta) \le \frac{\EE(|X|^p)}{\delta^p}.$$
\end{propositionbox}
\begin{proofbox}
    $$\EE(|X|^p) = \int |X|^p \dd \PP \ge \int_{|x| > \delta} |X|^p \dd \PP \ge \delta^p \cdot \PP(|X| > \delta).$$
\end{proofbox}
\begin{remarkbox}
    The case $p=1$ is \emph{Markov's inequality}: for any
    nonnegative random variable $Z$ and any $\delta>0$,
    \begin{equation}\label{eq:markov}
        \PP(Z>\delta)\leq\frac{\EE(Z)}{\delta}.
    \end{equation}
    Indeed, $Z\geq\delta\mathbbm{1}_{\{Z>\delta\}}$, so taking
    expectations gives $\EE(Z)\geq\delta\PP(Z>\delta)$.
    Thus a small expectation bounds the probability of exceeding
    any fixed positive threshold.
    Applying the proposition with $p=2$ to $X-\EE(X)$ gives
    \emph{Chebyshev's inequality}:
    \begin{equation}\label{eq:chebyshev}
        \PP(|X-\EE(X)|>\delta)\leq\frac{\Var(X)}{\delta^2}.
    \end{equation}
\end{remarkbox}
